\documentclass[12pt,a4paper]{article}

\usepackage[a4paper,margin=18mm]{geometry}
\usepackage[T2A]{fontenc}
\usepackage[utf8]{inputenc}
\usepackage[russian]{babel}
\usepackage{amsmath,amssymb,mathtools}
\usepackage{array,booktabs}
\usepackage{xcolor}
\usepackage{hyperref}
\usepackage{tikz}

\hypersetup{
  colorlinks=true,
  linkcolor=blue!55!black,
  urlcolor=blue!55!black
}

\setlength{\parindent}{0pt}
\setlength{\parskip}{0.55em}
\setlength{\tabcolsep}{3pt}
\renewcommand{\arraystretch}{1.25}

\begin{document}

\begin{titlepage}
\centering
\vspace*{\fill}

{\LARGE\bfseries Ревью домашней работы\par}

\vspace{1.2cm}

{\large Автор: Лёвин Артём Александрович\par}

\vspace{0.6cm}

{\large 12 сентября 2026 года\par}

\vspace*{\fill}

\begin{tikzpicture}
\draw[line width=0.7pt] (0,0) .. controls (2,-0.25) and (4,0.25) .. (6,0);
\end{tikzpicture}

\end{titlepage}

\newpage

\section*{Сводная таблица}

{\small
\begin{center}
\begin{tabular}{p{0.11\linewidth} p{0.18\linewidth} p{0.28\linewidth} p{0.14\linewidth} p{0.16\linewidth}}
\toprule
\textbf{№ задания} & \textbf{Тема} & \textbf{Суть ошибки} & \textbf{Ваш ответ} & \textbf{Правильный ответ} \\
\midrule
\hyperref[task:2026-09-11-2]{2 (11.09)} &
Равномерное движение двух тел &
Не найдена координата точки встречи: решение остановлено после нахождения времени &
$t=4\ \text{с}$ &
$x_{\text{встр}}=12\ \text{м}$ при $t=4\ \text{с}$ \\
\bottomrule
\end{tabular}
\end{center}
}

\newpage
\vspace*{\fill}
\begin{center}
{\LARGE\bfseries Разбор ошибок}
\end{center}
\vspace*{\fill}
\newpage

\phantomsection\label{task:2026-09-11-2}
\section*{Задание 2 от 11.09.2026}

\subsection*{Текст задания}
Начальные координаты двух тел:
\[
x_{01}=24\ \text{м},\qquad x_{02}=-4\ \text{м}.
\]
Скорости:
\[
v_{x1}=-3\ \text{м/с},\qquad v_{x2}=4\ \text{м/с}.
\]
Составьте законы движения и найдите точку встречи.

\subsection*{Ваше решение}
В записи правильно составлены законы движения:
\[
x_1=24-3t,
\qquad
x_2=-4+4t.
\]
Далее координаты тел при встрече приравнены:
\[
24-3t=-4+4t,
\]
откуда получено
\[
28=7t,
\qquad
 t=4\ \text{с}.
\]

\subsection*{Ваш ответ}
\[
t=4\ \text{с}.
\]

\subsection*{Ошибка}
\textcolor{red}{Ошибка:} решение не доведено до требуемой координаты точки встречи.

\subsection*{Где именно возникает ошибка}
Последний полностью корректный шаг:
\[
t=4\ \text{с}.
\]
Здесь вычислительной ошибки нет. Неполнота возникает на следующем этапе: после нахождения времени решение заканчивается, хотя по условию требуется найти \emph{точку встречи}, то есть координату, в которой оба тела окажутся одновременно.

Иными словами, строка
\[
\boxed{t=4\ \text{с}}
\]
является необходимым промежуточным результатом, но ещё не завершает задание.

\subsection*{Почему это неверно}
Для равномерного движения координата тела задаётся законом
\[
x=x_0+v_xt.
\]
Время встречи определяется из условия равенства координат:
\[
x_1(t)=x_2(t).
\]
Это даёт момент, когда тела находятся в одном месте. Однако само место встречи определяется отдельным шагом: найденное время необходимо подставить в закон движения любого из тел.

В данном случае найдено
\[
t=4\ \text{с}.
\]
Чтобы выполнить требование задачи полностью, нужно вычислить
\[
x_{\text{встр}}=x_1(4)
\]
или
\[
x_{\text{встр}}=x_2(4).
\]
Оба вычисления обязаны дать одну и ту же координату. Это одновременно служит проверкой решения.

\textcolor{blue!60!black}{Ключевая идея:} при задаче о встрече полезно разделять два результата: \emph{когда} произошла встреча и \emph{где} она произошла. Равенство законов движения сначала даёт время, а подстановка этого времени в любой закон движения даёт координату встречи.

\subsection*{Исправление от последнего правильного шага}
Уже найдено:
\[
t=4\ \text{с}.
\]
Подставим это время в первый закон движения:
\[
x_{\text{встр}}=24-3\cdot 4.
\]
Вычисляем:
\[
x_{\text{встр}}=24-12=12\ \text{м}.
\]
Следовательно, недостающий заключительный шаг:
\[
\boxed{x_{\text{встр}}=12\ \text{м}}.
\]

\subsection*{Полное правильное решение}
Для первого тела используем формулу равномерного движения
\[
x=x_0+v_xt.
\]
Так как
\[
x_{01}=24\ \text{м},\qquad v_{x1}=-3\ \text{м/с},
\]
получаем
\[
x_1=24-3t.
\]

Для второго тела:
\[
x_{02}=-4\ \text{м},\qquad v_{x2}=4\ \text{м/с},
\]
поэтому
\[
x_2=-4+4t.
\]

В момент встречи координаты равны:
\[
x_1=x_2.
\]
Значит,
\[
24-3t=-4+4t.
\]
Перенесём слагаемые с $t$ вправо, а числа влево:
\[
24+4=4t+3t.
\]
Получаем
\[
28=7t.
\]
Отсюда
\[
t=4\ \text{с}.
\]

Теперь найдём координату встречи. Подставим $t=4$ в первый закон движения:
\[
x_{\text{встр}}=24-3\cdot4=24-12=12\ \text{м}.
\]

Таким образом, тела встречаются через
\[
4\ \text{с}
\]
в точке с координатой
\[
12\ \text{м}.
\]

\subsection*{Проверка}
Проверим координату по закону движения второго тела:
\[
x_2(4)=-4+4\cdot4=-4+16=12\ \text{м}.
\]
Получено то же значение:
\[
x_1(4)=x_2(4)=12\ \text{м}.
\]
Значит, момент и координата встречи найдены согласованно.

\subsection*{Итог}
\textcolor{teal!60!black}{Итог:} составление законов движения и вычисление времени выполнены верно. Для полного ответа не хватило последнего шага --- подстановки $t=4\ \text{с}$ в один из законов движения. Полный результат:
\[
\boxed{t_{\text{встр}}=4\ \text{с},\qquad x_{\text{встр}}=12\ \text{м}}.
\]

\subsection*{Задача на закрепление}
Начальные координаты двух тел равны
\[
x_{01}=17\ \text{м},\qquad x_{02}=-7\ \text{м},
\]
а их скорости равны
\[
v_{x1}=-2\ \text{м/с},\qquad v_{x2}=4\ \text{м/с}.
\]
Составьте законы движения и найдите время и координату встречи.

\vfill
\begin{center}
\small Лёвин Артём Александрович эксклюзивно для Володи Хачатуряна
\end{center}

\end{document}
